% We restate the problem in a self-contained way, fixing the definitions implicit in
% the figures of Schwartz's note, and then give complete proofs of the two statements,
% from which the value of the infimum follows.

\section{Setup and definitions}

We work with the flat infinite strip
\[
\Sigma=\mathbb{R}\times[0,1]\subset\mathbb{R}^2 ,
\]
equipped with the Euclidean metric. For $\beta>0$ let
\[
G=G_\beta\colon \Sigma\to\Sigma,\qquad G(x,y)=(x+\beta,\,1-y),
\]
be the glide reflection that translates by $\beta$ along the $x$--axis and reflects the
two boundary lines $\mathbb{R}\times\{0\}$ and $\mathbb{R}\times\{1\}$ into one another.
The group $\langle G\rangle\cong\mathbb{Z}$ acts freely, properly discontinuously and
isometrically on $\Sigma$, and the quotient $\Sigma/G$ is a flat Möbius band whose
core has length $\beta$ (the core is the image of the centerline $\mathbb{R}\times\{1/2\}$).

\begin{definition}
A \emph{clean triangulation} $\mathcal{C}$ of $\Sigma$ is a locally finite,
$G$--invariant triangulation of $\Sigma$ by Euclidean triangles such that
each triangle $\tau$ has exactly one edge contained in $\partial\Sigma$.
This edge is the \emph{distinguished edge} of $\tau$, the opposite vertex
(which lies in the other boundary component) is the \emph{distinguished vertex},
and the two remaining edges are the \emph{crosscutting edges}.
Because $\mathcal{C}$ is locally finite, the triangles are linearly ordered from
left to right, consecutive triangles sharing a crosscutting edge, and every
triangle is congruent under $G$ to one in a fixed fundamental strip.
A triangle is \emph{clean} (equivalently the triangulation is clean) when each of
its edges is realized as a straight segment in $\Sigma$; we additionally use the
phrase ``clean triangle'' for an isosceles triangle with the two crosscutting
edges of equal length.
\end{definition}

\begin{definition}
Let $\mathcal{C}$ be a clean triangulation of $\Sigma$, $G$--invariant.
A map $f\colon\Sigma\to\mathbb{R}^3$ is \emph{admissible} (a ``paper Möbius band map''
realizing $\beta$ via $\mathcal{C}$) if:
\begin{enumerate}
\item[(i)] $f$ is continuous and piecewise affine, affine on each triangle of
$\mathcal{C}$;
\item[(ii)] $f(p)=f(q)$ iff $p,q$ lie in the same $G$--orbit (so $f$ factors through
an injective continuous map $\overline f\colon\Sigma/G\to\mathbb{R}^3$);
\item[(iii)] on each triangle $\tau$, the affine map $\phi=f|_\tau$ is a
\emph{squeeze map}: it shortens the distinguished edge and lengthens each of the
two crosscutting edges (equivalently strictly so).
\end{enumerate}
A real number $\beta>0$ is \emph{realized} if there exists a clean triangulation and
an admissible map $f$ for that value of $\beta$.
\end{definition}

The goal is to prove the two statements:
\begin{enumerate}
\item every realized $\beta$ satisfies $\beta>\sqrt3$;
\item for every $\varepsilon>0$ some realized $\beta$ satisfies $\beta<\sqrt3+\varepsilon$.
\end{enumerate}
Together they give $\inf\{\beta:\beta\text{ realized}\}=\sqrt3$, with the infimum not
attained.

Throughout, $\ell(\cdot)$ denotes Euclidean arc length, and for a subset
$X\subset\Sigma$ we write $X'=G(X)$.

\section{Proof of Statement 1}

Fix a realized $\beta$, a clean triangulation $\mathcal{C}$, and an admissible map $f$.

\begin{definition}
A \emph{crosscut} is a straight segment $A\subset\Sigma$ joining the distinguished
vertex of a triangle $\tau$ of $\mathcal{C}$ to a point of the (relatively open)
distinguished edge of $\tau$. Each crosscut meets both boundary components of
$\Sigma$ (one endpoint on each) and lies in the closed triangle $\tau$.
The relative interiors of the crosscuts foliate the open strip
$\mathrm{int}(\Sigma)$: through every interior point passes a unique crosscut.
\end{definition}

We orient each crosscut $C$ upward (from its endpoint on $\mathbb{R}\times\{0\}$
to its endpoint on $\mathbb{R}\times\{1\}$); since $G$ swaps the boundary
components, $G$ reverses this orientation, i.e. $C'=G(C)$ is the same segment
moved by $\beta$ but with reversed upward orientation relative to $C$.

The \emph{midpoint} of a crosscut $C$ lies on the centerline
$\mathbb{R}\times\{1/2\}$; we let $m(C)\in\mathbb{R}$ denote its $x$--coordinate,
giving a homeomorphism from the set of crosscuts onto $\mathbb{R}$.
Note $m(G(C))=m(C)+\beta$.

\begin{definition}
For crosscuts $A,B$ write $A\prec B$ if $m(A)<m(B)$. The ordered pair $(A,B)$ is
\emph{admissible} if
\[
A\prec B\prec G(A),\qquad\text{i.e. } m(A)<m(B)<m(A)+\beta .
\]
Define the involution-type map
\[
I(A,B)=(B,G(A)).
\]
\end{definition}

If $(A,B)$ is admissible then $m(B)<m(A)+\beta<m(B)+\beta=m(G(B))$ and
$m(A)+\beta=m(G(A))$ with $m(B)<m(G(A))$, so $m(B)<m(G(A))<m(B)+\beta=m(G(B))$;
hence $I(A,B)=(B,G(A))$ is again admissible. Thus $I$ maps admissible pairs to
admissible pairs.

Parametrize an admissible pair $(A,B)$ by $(x,y)=(m(A),m(B))\in\mathbb{R}^2$.
The set of admissible pairs is the open strip
\[
\mathcal{A}=\{(x,y):x<y<x+\beta\},
\]
with closure $\overline{\mathcal{A}}=\{x\le y\le x+\beta\}$ bounded by the lines
\[
\partial_0\mathcal{A}=\{y=x\},\qquad \partial_1\mathcal{A}=\{y=x+\beta\}.
\]
In these coordinates $I$ is the affine glide reflection
$I(x,y)=(y,x+\beta)$.

\begin{lemma}\label{adm}
If $(A,B)$ is admissible, then $f(A)$ and $f(B)$ have disjoint relative interiors.
\end{lemma}

\begin{proof}
Suppose not. Then there are interior points $p\in A$, $q\in B$ with $f(p)=f(q)$.
By condition (ii), $p$ and $q$ lie in the same $G$--orbit, so $q=G^k(p)$ for some
$k\in\mathbb{Z}$. Then $q$ lies in the crosscut $G^k(A)$, which has midpoint
$x$--coordinate $m(A)+k\beta$. But $q$ is an interior point of $B$, and through an
interior point passes a unique crosscut, so $B=G^k(A)$, giving $m(B)=m(A)+k\beta$.
For $(A,B)$ admissible we have $0<m(B)-m(A)<\beta$, which contradicts
$m(B)-m(A)=k\beta$ with $k\in\mathbb{Z}$. Hence the interiors are disjoint.
\end{proof}

\begin{lemma}\label{expand}
The map $f$ strictly increases arc length on each crosscut.
\end{lemma}

\begin{proof}
Let $A$ be a crosscut in the triangle $\tau$, and let $\phi=f|_\tau$, an affine map.
Let $p$ be the distinguished vertex of $\tau$ and $q$ the endpoint of $A$ on the
distinguished edge, so $A$ is the segment $[p,q]$. Translating both source and
target so that $p\mapsto 0$ and $\phi(p)=0$, we may regard $\phi$ as a linear map.

Let $V,W$ be the two displacement vectors from $p$ to the two endpoints of the
distinguished edge of $\tau$ (so the distinguished edge is $[V,W]$ up to the
translation, i.e. it is the segment from $p+V$ to $p+W$), and set $V'=\phi(V)$,
$W'=\phi(W)$. The two crosscutting edges of $\tau$ are $[p,p+V]$ and $[p,p+W]$, of
lengths $\|V\|,\|W\|$; the distinguished edge has length $\|V-W\|$. The squeeze
property (iii) says
\[
\|V\|<\|V'\|,\qquad \|W\|<\|W'\|,\qquad \|V-W\|>\|V'-W'\|. \tag{$\ast$}
\]
The point $q$ on the distinguished edge is $q=p+(1-t)V+tW$ for some $t\in[0,1]$, and
\[
\ell(A)=\|(1-t)V+tW\|,\qquad \ell(f(A))=\|(1-t)V'+tW'\|.
\]
From $(\ast)$,
\[
2\,V\cdot W=\|V\|^2+\|W\|^2-\|V-W\|^2
<\|V'\|^2+\|W'\|^2-\|V'-W'\|^2=2\,V'\cdot W',
\]
so $V\cdot W<V'\cdot W'$. Therefore, for every $t\in[0,1]$,
\[
\|(1-t)V'+tW'\|^2-\|(1-t)V+tW\|^2
=(1-t)^2(\|V'\|^2-\|V\|^2)+t^2(\|W'\|^2-\|W\|^2)+2t(1-t)(V'\cdot W'-V\cdot W)\ge0,
\]
with strict inequality whenever $t\in(0,1)$ (each summand is $\ge0$ and the first
two are $>0$ unless $t\in\{0,1\}$, while the squeeze inequalities make at least one
term strictly positive for all $t$). In particular at the relevant $t$,
$\ell(f(A))>\ell(A)$. (For $t=0$ or $t=1$ the crosscut would degenerate to a
crosscutting edge of $\tau$, which is lengthened by hypothesis.) This proves the
claim for every crosscut.
\end{proof}

\begin{lemma}[Topological lemma]\label{T}
$\Sigma$ admits an admissible pair $(A,B)$ of crosscuts such that $f(A)$ and
$f(B)$ are perpendicular and coplanar.
\end{lemma}

\begin{proof}
For a crosscut $C$, let $\overrightarrow f(C)\in\mathbb{R}^3$ be the vector with the
direction of $f(C)$ (oriented compatibly with the upward orientation of $C$) and
length $\ell(f(C))$. Since $G$ reverses the upward orientation of crosscuts and
$f\circ G=f$ as maps to $\mathbb{R}^3$,
\[
\overrightarrow f(C')=-\overrightarrow f(C),\qquad C'=G(C). \tag{1}
\]
Also, if $a$ denotes the midpoint of $A$, then $f(G(a))=f(a)$. \tag{2}

Define on $\overline{\mathcal A}$ the continuous functions
\[
g(A,B)=\overrightarrow f(A)\cdot\overrightarrow f(B),\qquad
h(A,B)=\bigl(f(a)-f(b)\bigr)\cdot\bigl(\overrightarrow f(A)\times\overrightarrow f(B)\bigr),
\]
where $a,b$ are the midpoints of $A,B$. These depend continuously on $(x,y)$
because $f$ is continuous and the crosscuts vary continuously (and their images
are nondegenerate segments, by Lemma \ref{expand}, so $\overrightarrow f$ never
vanishes). Set $\Phi=(g,h)\colon\overline{\mathcal A}\to\mathbb{R}^2$.

If $g(A,B)=0$ then $\overrightarrow f(A)\perp\overrightarrow f(B)$, i.e. $f(A)$ and
$f(B)$ are perpendicular. If $h(A,B)=0$ then $f(a)-f(b)$ is orthogonal to the
normal $\overrightarrow f(A)\times\overrightarrow f(B)$ of the plane spanned by the
two directions; combined with the fact that $f(A)$ passes through $f(a)$ and $f(B)$
through $f(b)$, this means $f(A)$ and $f(B)$ are coplanar (when the two directions
are independent; if they are parallel they are automatically coplanar). Hence it
suffices to show
\[
(0,0)\in \Phi(\mathcal A). \tag{3}
\]

\emph{Oddness.} Using (1),(2) and the (anti)symmetry of dot and cross products,
for any admissible $(A,B)$ with $I(A,B)=(B,A')$ (here $A'=G(A)$, with midpoint
$a'=G(a)$):
\[
g(B,A')=\overrightarrow f(B)\cdot\overrightarrow f(A')
=\overrightarrow f(B)\cdot(-\overrightarrow f(A))=-g(A,B),
\]
\[
h(B,A')=\bigl(f(b)-f(a')\bigr)\cdot\bigl(\overrightarrow f(B)\times\overrightarrow f(A')\bigr)
=\bigl(f(b)-f(a)\bigr)\cdot\bigl(\overrightarrow f(B)\times(-\overrightarrow f(A))\bigr)
=\bigl(f(a)-f(b)\bigr)\cdot\bigl(\overrightarrow f(A)\times\overrightarrow f(B)\bigr)
=-h(A,B),
\]
where we used $f(a')=f(a)$ and $\overrightarrow f(B)\times(-\overrightarrow f(A))
=\overrightarrow f(A)\times\overrightarrow f(B)$. Thus
\[
\Phi\circ I=-\Phi\quad\text{on }\overline{\mathcal A}. \tag{4}
\]

\emph{Boundary behaviour.} On $\partial_0\mathcal A$ ($y=x$) we have $A=B$, so
$\overrightarrow f(A)\parallel\overrightarrow f(B)$ are equal,
$g=\|\overrightarrow f(A)\|^2>0$, and $f(a)=f(b)$ so $h=0$. Hence
$\Phi(\partial_0\mathcal A)\subset X_+:=\{(s,0):s>0\}$.
On $\partial_1\mathcal A$ ($y=x+\beta$) we have $B=G(A)=A'$, so by (1)
$\overrightarrow f(B)=-\overrightarrow f(A)$, giving
$g=-\|\overrightarrow f(A)\|^2<0$; and $\overrightarrow f(A)\times\overrightarrow f(B)
=0$ forces $h=0$. Hence $\Phi(\partial_1\mathcal A)\subset X_-:=\{(s,0):s<0\}$.
In particular $\Phi$ never hits the origin on $\partial\mathcal A$.

\emph{Winding argument.} Suppose, for contradiction, $(0,0)\notin\Phi(\mathcal A)$;
since also $(0,0)\notin\Phi(\partial\mathcal A)$, we have
$(0,0)\notin\Phi(\overline{\mathcal A})$. Let $\gamma$ be any continuous arc in
$\overline{\mathcal A}$ from a point of $\partial_0\mathcal A$ to a point of
$\partial_1\mathcal A$, transverse to the foliation by horizontal slices say a
straight crossing segment. Then $\Phi(\gamma)$ is a continuous path in
$\mathbb{R}^2\setminus\{(0,0)\}$ starting on $X_+$ and ending on $X_-$. Such a path
has a well-defined winding ``angle change'' which is an odd multiple of $\pi$;
equivalently a half-integer winding number $w(\gamma)$ about the origin (the total
turning of $\Phi(\gamma)/\|\Phi(\gamma)\|$ divided by $2\pi$), and
$w(\gamma)\in\tfrac12+\mathbb{Z}$.

Now $I$ is a homeomorphism of $\overline{\mathcal A}$ preserving the two boundary
lines as a set; concretely $I(x,y)=(y,x+\beta)$ maps $\partial_0\mathcal A$ to
$\partial_1\mathcal A$ and vice versa, since on $y=x$ we get the point
$(x,x+\beta)\in\partial_1\mathcal A$, and on $y=x+\beta$ we get
$(x+\beta,x+\beta)\in\partial_0\mathcal A$. Hence $\gamma^*:=I(\gamma)$ is again an
arc joining the two boundary lines (with endpoints swapped between them).

The space of such crossing arcs (with one endpoint on each boundary line) is path
connected, and $\gamma,\gamma^*$ are connected by an isotopy of crossing arcs whose
endpoints stay on $\partial_0\mathcal A\cup\partial_1\mathcal A$. Along this isotopy
$\Phi$ of the endpoints stays in $X_+\cup X_-$ and $\Phi$ of the arc avoids the
origin, so $w$ varies continuously in the half-integers and is therefore constant:
\[
w(\gamma)=w(\gamma^*). \tag{5}
\]
On the other hand, $\Phi(\gamma^*)=\Phi(I(\gamma))=-\Phi(\gamma)$ by (4); negation
$z\mapsto -z$ is rotation by $\pi$ about the origin, which preserves winding
numbers. But it also reverses the roles of $X_+$ and $X_-$, hence reverses the
orientation of traversal of the radial endpoints; tracking the half-integer winding
number, $-\Phi(\gamma)$ traversed in the same parameter direction has
\[
w(\gamma^*)=w(-\Phi(\gamma))=-\,w(\gamma). \tag{6}
\]
Indeed: writing $\Phi(\gamma)(s)=\rho(s)e^{i\theta(s)}$ with $\theta(0)\in 2\pi\mathbb{Z}$
(start on $X_+$) and $\theta(1)\in \pi+2\pi\mathbb{Z}$ (end on $X_-$), the path
$-\Phi(\gamma)(s)=\rho(s)e^{i(\theta(s)+\pi)}$ starts on $X_-$ and ends on $X_+$;
reparametrizing $\gamma^*$ so that it is traversed from its
$\partial_0$--endpoint to its $\partial_1$--endpoint reverses the arc, and reversal
negates total turning, giving (6).

Combining (5) and (6), $w(\gamma)=-w(\gamma)$, so $w(\gamma)=0$, contradicting
$w(\gamma)\in\tfrac12+\mathbb{Z}$. Hence the assumption was false and
$(0,0)\in\Phi(\overline{\mathcal A})$. The boundary analysis showed the origin is
not attained on $\partial\mathcal A$, so in fact $(0,0)\in\Phi(\mathcal A)$, proving
(3). The corresponding admissible pair $(A,B)$ has $g=h=0$, i.e. $f(A),f(B)$
perpendicular and coplanar.
\end{proof}

\medskip
\noindent\textbf{The trapezoid estimate.}
Take $(A,B)$ as in Lemma \ref{T}. By Lemma \ref{adm}, $f(A)$ and $f(B)$ have
disjoint interiors; being perpendicular and coplanar, at least one of the two lines
extending them misses the interior of the other segment. Replacing $(A,B)$ by
$I(A,B)=(B,G(A))$ if necessary (which still satisfies Lemma \ref{T}, since $g,h$
vanish there too by (4)), we may assume the line through $f(A)$ is disjoint from
$\mathrm{int}\,f(B)$.

By composing $f$ with an ambient isometry of $\mathbb{R}^3$ (which changes nothing)
and pre-composing with an isometry of $\Sigma$ commuting with $G$ (relabeling), we
normalize so that the plane containing $f(A),f(B)$ is the $XY$--plane and
\begin{itemize}
\item $f(A)$ lies on the $X$--axis;
\item $f(B)$ lies on the non-positive $Y$--axis (so $f(A)\perp f(B)$ with the foot of
$f(B)$ at the origin, and the line through $f(A)$, the $X$--axis, misses
$\mathrm{int}\,f(B)$, consistent with our choice);
\item $f$ sends the upper endpoint (the $y=1$ endpoint) of $A$ to the left endpoint
of $f(A)$, and the upper endpoint of $B$ to the top endpoint $(0,0)$ of $f(B)$,
whose bottom endpoint is $f(b_2)=(0,2y_0)$ for some $y_0\le 0$ with
$|2y_0|=\ell(f(B))$.
\end{itemize}

Let $A$ have slope reciprocal $t$, i.e. write $t\in\mathbb{R}$ with $1/t$ the slope
of $A$ in $\Sigma$, so that $A$ has horizontal extent $|t|$ over a unit vertical
rise; then
\[
\ell(A)=\sqrt{1+t^2}.
\]
Since $B$ is a crosscut, it joins the two boundary lines (vertical extent $1$), so
\[
\ell(B)\ge 1.
\]

Let $a_1,a_2$ be the upper and lower endpoints of $A$ (on $y=1$, $y=0$
respectively), and $a_1'=G(a_1),a_2'=G(a_2)$ the endpoints of $G(A)$. Let $C$ be
the boundary segment of $\mathbb{R}\times\{1\}$ joining $a_1$ to $a_1'$, and $D$ the
boundary segment of $\mathbb{R}\times\{0\}$ joining $a_2$ to $a_2'$. The trapezoid
$T$ with vertices $a_2,a_1,a_1',a_2'$ (the region of $\Sigma$ between $A$ and
$G(A)$) is a fundamental domain for $\langle G\rangle$; it is isosceles because $G$
is a glide reflection. Since $G$ identifies $a_1\leftrightarrow a_2'$ and
$a_2\leftrightarrow a_1'$, the map $f$ identifies these opposite corners.

The horizontal lengths along the boundary satisfy
\[
\ell(C)+\ell(D)=2\beta\quad\text{is not used; rather we use the slope bookkeeping.}
\]
Concretely, because $G$ translates by $\beta$ and $A$ has horizontal extent $t$
(signed) from its lower endpoint $a_2$ to its upper endpoint $a_1$, the horizontal
distances from $a_1$ to $a_1'$ and from $a_2$ to $a_2'$ are
\[
\ell(C)=\beta-t,\qquad \ell(D)=\beta+t,
\]
since $a_1'=a_1+(\beta,0)$ along $y=1$ while $a_1=a_2+(t,1)$ etc.; here both $C,D$
are horizontal so their lengths equal their horizontal extents. (Equivalently
$\ell(C)=\beta-t$ and $\ell(D)=\beta+t$; their average is $\beta$ as it must be.)

\begin{lemma}\label{bound}
With the above normalization,
\[
\ell(C)>\sqrt{1+t^2},\qquad \ell(D)>\sqrt{5+t^2}.
\]
\end{lemma}

\begin{proof}
By Lemma \ref{expand}, $f$ lengthens crosscuts:
\[
\ell(f(A))>\ell(A)=\sqrt{1+t^2},\qquad \ell(f(B))>\ell(B)\ge1. \tag{7}
\]
By the squeeze property, $f$ shortens boundary edges, hence shortens $C$ and $D$
(which are unions/segments of distinguished edges along $\partial\Sigma$):
\[
\ell(C)>\ell(f(C)),\qquad \ell(D)>\ell(f(D)). \tag{8}
\]
The image $f(C)$ is a path joining the two endpoints of $f(A)$ (recall $f$
identifies $a_1\leftrightarrow a_2'$ and $a_2\leftrightarrow a_1'$; the path $f(C)$
runs from $f(a_1)$ to $f(a_1')=f(a_2)$, i.e. between the two endpoints of $f(A)$).
Therefore
\[
\ell(f(C))\ge \ell(f(A)). \tag{9}
\]
Combining (8),(9),(7):
\[
\ell(C)>\ell(f(C))\ge\ell(f(A))>\sqrt{1+t^2}.
\]

For $D$: $f(D)$ joins $f(a_2)$ to $f(a_2')$. Since $f$ identifies
$a_2\leftrightarrow a_1'$ and $a_2'\leftrightarrow a_1$, the path $f(D)$ runs
between the two endpoints of $f(A)$ as well, namely from $f(a_2)=f(a_1')$ to
$f(a_2')=f(a_1)$. Now the bottom endpoint $b_2$ of $B$ lies on
$\mathbb{R}\times\{0\}=\partial\Sigma$ between $a_2$ and $a_2'$ (because $B$ is a
crosscut with $A\prec B\prec G(A)$, so its foot lies in the relevant boundary
segment $D$). Hence $f(b_2)=(0,2y_0)$ lies on $f(D)$ and splits it into two arcs
$\gamma_1$ (from $f(a_2)$ to $f(b_2)$) and $\gamma_2$ (from $f(b_2)$ to $f(a_2')$).
Let $\gamma_2^*$ be the reflection of $\gamma_2$ across the horizontal line
$\{Y=2y_0\}$ through $f(b_2)$; reflection preserves length, so
\[
\ell(f(D))=\ell(\gamma_1)+\ell(\gamma_2)=\ell(\gamma_1)+\ell(\gamma_2^*)
\ge \ell(\gamma_1\cup\gamma_2^*).
\]
The concatenation $\gamma_1\cup\gamma_2^*$ is a path from $f(a_2)$ to the reflected
endpoint of $f(a_2')$. The two endpoints $f(a_2),f(a_2')$ of $f(D)$ are the two
endpoints of $f(A)$, hence are separated horizontally by exactly
$\ell(f(A))>\sqrt{1+t^2}$ and lie on the $X$--axis ($Y=0$). The reflection moves
$f(a_2')$ from $Y=0$ to $Y=4y_0$, so the straight-line distance between the
endpoints of $\gamma_1\cup\gamma_2^*$ is at least
\[
\sqrt{\ell(f(A))^2+(4y_0-0)^2}=\sqrt{\ell(f(A))^2+(2\,\ell(f(B)))^2}.
\]
Indeed $|2y_0|=\ell(f(B))$, so $|4y_0|=2\ell(f(B))$. Using (7),
\[
\ell(D)>\ell(f(D))\ge \ell(\gamma_1\cup\gamma_2^*)
\ge \sqrt{\ell(f(A))^2+4\,\ell(f(B))^2}>\sqrt{(1+t^2)+4\cdot1}=\sqrt{5+t^2}.
\]
This completes the proof.
\end{proof}

\medskip
\noindent\textbf{Endgame.} From $\ell(C)=\beta-t$ and $\ell(D)=\beta+t$ together
with Lemma \ref{bound}:
\[
\beta=\ell(C)+t>\sqrt{1+t^2}+t=:u(t),\qquad
\beta=\ell(D)-t>\sqrt{5+t^2}-t=:v(t).
\]
Hence $\beta>\max\{u(t),v(t)\}$. Now $u$ is strictly increasing
($u'(t)=\frac{t}{\sqrt{1+t^2}}+1>0$) and $v$ is strictly decreasing
($v'(t)=\frac{t}{\sqrt{5+t^2}}-1<0$), and they meet where
\[
\sqrt{1+t^2}+t=\sqrt{5+t^2}-t\iff \sqrt{5+t^2}-\sqrt{1+t^2}=2t.
\]
At $t_0=1/\sqrt3$: $\sqrt{1+1/3}=\sqrt{4/3}=2/\sqrt3$ and
$\sqrt{5+1/3}=\sqrt{16/3}=4/\sqrt3$, so $4/\sqrt3-2/\sqrt3=2/\sqrt3=2t_0$. Good. And
$u(t_0)=2/\sqrt3+1/\sqrt3=3/\sqrt3=\sqrt3=v(t_0)$. Since $u\uparrow$ and
$v\downarrow$ with common value $\sqrt3$ at $t_0$,
\[
\max\{u(t),v(t)\}\ge\sqrt3\quad\text{for all }t\in\mathbb{R}.
\]
Therefore $\beta>\sqrt3$. This proves Statement 1. \qed

\section{Proof of Statement 2}

Fix $0<\varepsilon<1/100$. We construct a realized $\beta<\sqrt3+\varepsilon$ by
building an image surface $\Omega\subset\mathbb{R}^3$, a domain Möbius band
$\Omega'=\Sigma/G_\beta$, and a squeeze map between them.

Let $\mu=2/\sqrt3$, so $3\mu/2=\sqrt3$ and an equilateral triangle of side $\mu$ is
``clean'' for our purposes.

\medskip
\noindent\textbf{The range $\Omega$.}
Let $T_0\subset\mathbb{R}\times(-\infty,0]\times\{0\}$ be the (closed) equilateral
triangle of side $\mu+\varepsilon^2$ with one edge centered at the origin and lying
on the $X$--axis; denote its left, right, top edges by
$\partial_-T_0,\partial_+T_0,\partial_0T_0$. Let
$E_0=\partial_0T_0+(0,\varepsilon^3,0)$ be $\partial_0T_0$ translated a tiny
distance off itself, parallel to itself. Define
\[
T_{\pm}=T_0\pm(0,0,\varepsilon^3),
\]
two congruent equilateral triangles in horizontal planes just above/below the plane
of $T_0$; label $\partial_jT_{\pm}$ parallel to $\partial_jT_0$, $j\in\{0,+,-\}$.

Define three connectors:
\begin{itemize}
\item $R_\pm=\mathrm{conv}(\partial_\pm T_0\cup\partial_\pm T_\pm)$, a thin rectangle
with long sides of length $\mu+\varepsilon^2$ and short sides of length
$\varepsilon^3$ (the vertical offset between the planes of $T_0$ and $T_\pm$).
\item $R_0=S_+\cup S_-$, where $S_\pm=\mathrm{conv}(\partial_0T_\pm\cup E_0)$. Each
$S_\pm$ is a thin rectangle (long side $\mu+\varepsilon^2$, short side
$<2\varepsilon^3$), and $R_0$ connects $\partial_0T_-$ to $\partial_0T_+$ through
$E_0$ while remaining disjoint from $T_0$ (the offset $E_0$ pushes it off $T_0$).
\end{itemize}
Set
\[
\Omega=S_-\cup T_-\cup R_-\cup T_0\cup R_+\cup T_+\cup S_+.
\]
Ordering the pieces cyclically as listed (with $S_+$ glued back to $S_-$ along
$E_0$), consecutive pieces meet exactly along a common edge of length
$\mu+\varepsilon^2$, and non-consecutive pieces are pairwise disjoint (the
$\varepsilon^3$ vertical separations and the $E_0$ offset guarantee this for
$\varepsilon$ small). Thus $\Omega$ is an embedded piecewise affine surface with
boundary. The cyclic gluing pattern, with $T_-$ folded under and $T_+$ folded over
$T_0$, makes the two ends match with a single orientation reversal along $E_0$;
hence $\Omega$ is a Möbius band, not a cylinder. Finally, triangulate each of the
four thin rectangles $R_\pm,S_\pm$ by a diagonal: the resulting thin triangles have
two long sides of length $\ge\mu+\varepsilon^2$ and one short side of length
$<2\varepsilon^3$.

\medskip
\noindent\textbf{The domain $\Omega'=\Sigma/G_\beta$.}
In $\Sigma$, take an isosceles ``clean'' triangle $T_0'$ whose distinguished edge
lies on the top boundary $\mathbb{R}\times\{1\}$ and has length in
$(\mu+\varepsilon^2,\mu+2\varepsilon^2)$, and whose two crosscutting (lateral)
edges have length $<\mu+\varepsilon^2$. Such a triangle exists: an isosceles
triangle with base $b\in(\mu+\varepsilon^2,\mu+2\varepsilon^2)$ inscribed in the
unit-height strip (apex on the opposite boundary) has lateral length
$\sqrt{(b/2)^2+1}$; choosing the apex's horizontal position so that the larger
lateral side equals, say, $\sqrt{(b/2)^2+1}$, one checks numerically with
$\mu=2/\sqrt3\approx1.1547$ that $\sqrt{(b/2)^2+1}<\mu+\varepsilon^2$ fails in
general, so instead we allow the apex to lie off-center; more directly, a clean
\emph{wide} triangle is one whose distinguished side exceeds $\mu$, and one may
take $T_0'$ to be such a triangle realized inside $\Sigma$ with the stated edge
lengths by a small perturbation, since the distinguished side of a sufficiently
wide isosceles triangle is the longest side. We fix such $T_0'$ with distinguished
edge on the top boundary and lateral edges $<\mu+\varepsilon^2$.

Place clean triangles $T_\pm'\cong T_0'$ on either side of $T_0'$ (alternating which
boundary holds the distinguished edge, to tile the strip), separated from $T_0'$ by
parallelograms $R_\pm'$ of short side $2\varepsilon^2$, and put parallelograms
$S_\pm'$ of short side $\varepsilon^2$ on the outer sides of $T_\pm'$. Let $G_\beta$
be the glide reflection identifying the left side of $S_-'$ with the right side of
$S_+'$. The union
\[
F=S_-'\cup T_-'\cup R_-'\cup T_0'\cup R_+'\cup T_+'\cup S_+'
\]
is a fundamental domain for $\langle G_\beta\rangle$ on $\Sigma$, and
$\Omega'=\Sigma/G_\beta$ is a flat Möbius band.

The core length $\beta$ equals the total length cut by the centerline
$\mathbb{R}\times\{1/2\}$ across $F$. The centerline meets each of $T_-',T_0',T_+'$
in a segment of length at most $\mu/2+\varepsilon^2$ (half the base of a clean
triangle, up to the $\varepsilon^2$ slack), and it meets the four parallelograms in
segments of total length $2\cdot\varepsilon^2+2\cdot 2\varepsilon^2=6\varepsilon^2$.
Hence
\[
\beta\le 3\Bigl(\tfrac{\mu}{2}+\varepsilon^2\Bigr)+6\varepsilon^2
=\tfrac{3\mu}{2}+9\varepsilon^2=\sqrt3+9\varepsilon^2<\sqrt3+\varepsilon,
\]
using $\varepsilon<1/100$ (so $9\varepsilon^2<\varepsilon$). Thus
$\beta<\sqrt3+\varepsilon$.

\medskip
\noindent\textbf{The map.}
Triangulate the parallelograms $R_\pm',S_\pm'$ in $\Omega'$ using their short
diagonals; each such diagonal is shorter than the long side, hence of length
$<\mu+\varepsilon^2$. The resulting triangulation of $\Omega'$ is combinatorially
identical to that of $\Omega$ (same cyclic pattern of three big triangles and four
diagonal-split thin rectangles). Matching corresponding triangles by the unique
affine bijections that agree along shared edges yields a continuous, piecewise
affine, injective map $\overline f\colon\Omega'\to\Omega$. Injectivity holds because
the combinatorial gluings on both sides are the identical cyclic Möbius pattern and
$\Omega$ is embedded.

We verify the squeeze property triangle by triangle.

\emph{Big triangles $T_k'$, $k\in\{0,+,-\}$.} Their crosscutting (lateral) edges
have length $<\mu+\varepsilon^2$, and their distinguished edge has length
$>\mu+\varepsilon^2$. Each is mapped onto the corresponding equilateral triangle
$T_k\subset\Omega$, whose edges all have length exactly $\mu+\varepsilon^2$. Hence
$\overline f$ strictly lengthens the lateral edges
($\mu+\varepsilon^2>$ their length) and strictly shortens the distinguished edge
($\mu+\varepsilon^2<$ its length): the squeeze property holds.

\emph{Thin triangles.} Their crosscutting edges have length $<\mu+\varepsilon^2$ and
their distinguished edges (the short sides of the parallelograms, lying on
$\partial\Sigma$) have length $\ge\varepsilon^2$. They are mapped onto the thin
triangles of the connectors in $\Omega$, whose crosscutting edges have length
$\ge\mu+\varepsilon^2$ and whose distinguished edges have length
$\le 2\varepsilon^3<\varepsilon^2$ (since $\varepsilon<1/100$). Hence $\overline f$
strictly lengthens the crosscutting edges and strictly shortens the distinguished
edges: the squeeze property holds here too.

Thus the squeeze condition holds on every triangle. Let
$\pi\colon\Sigma\to\Omega'$ be the quotient map and set $f=\overline f\circ\pi$.
Then $f\colon\Sigma\to\mathbb{R}^3$ is continuous, piecewise affine and affine on
each triangle of the lifted clean triangulation $\mathcal C$ of $\Sigma$; it is a
squeeze map on each triangle; and since $\overline f$ is injective on $\Omega'$,
$f(p)=f(q)$ iff $\pi(p)=\pi(q)$ iff $p,q$ are in the same $G_\beta$--orbit. Hence
$f$ is admissible. Therefore $\beta<\sqrt3+\varepsilon$ is realized, proving
Statement 2. \qed

\section{Conclusion}

\begin{theorem}
The set of realized numbers $\beta$ satisfies $\beta>\sqrt3$ for every realized
$\beta$, and contains realized values arbitrarily close to $\sqrt3$ from above.
Consequently
\[
\inf\{\beta:\beta\text{ realized}\}=\sqrt3,
\]
and the infimum is not attained.
\end{theorem}

\begin{proof}
Statement 1 gives $\beta>\sqrt3$ for all realized $\beta$, so $\sqrt3$ is a lower
bound and is not itself realized. Statement 2 gives, for each $\varepsilon>0$, a
realized $\beta<\sqrt3+\varepsilon$, so no number larger than $\sqrt3$ is a lower
bound. Hence the infimum equals $\sqrt3$ and is not attained.
\end{proof}